\documentclass[../differential_topology.tex]{subfiles}
\begin{document}
\section{Aula 03 - 20 de Agosto, 2026}
\subsection{Motivações}
\begin{itemize}
	\item O Plano Tangente
\end{itemize}
\subsection{O Plano Tangente}
\begin{def*}
	Sejam M uma variedade de dimensão m, p um ponto de M e \(\varphi :U\subseteq \mathbb{R}^{m}\rightarrow V\subseteq M\subseteq \mathbb{R}^{k}\) uma parametrização de M em p, onde V é aberto e contém p, com \(\varphi(x_{0}) = p\). Definimos o \textbf{plano tangente de M em p} como sendo
	\[
		T_{p}M = \varphi'(x_{0})\cdot \mathbb{R}^{m}
	\]
\end{def*}
Em particular, o plano tangente é bem-definido e tem dimensão igual à da variedade; podemos supor que \(\varphi \) é gráfico de uma função!
\begin{def*}
	Dada \(M\subseteq \mathbb{R}^{m}\) uma variedade suave de dimensão m, definimos o \textbf{fibrado tangente de M} como
	\[
		TM = \{(p, v): p\in M\; \& \; v\in T_{p}M\} = \bigcup_{p\in M}^{}(\{\} \times T_{p}M)\subseteq \mathbb{R}^{k}\times \mathbb{R}^{k}.\; \square
	\]
\end{def*}
Essa é uma variedade de dimensão 2m, e podemos conseguir uma parametrização para ele considerando uma parametrização de M \(\varphi :U\subseteq \mathbb{R}^{m}\rightarrow V\subseteq M\), colocando depois
\begin{align*}
	\Phi : & U\times \mathbb{R}^{m}\rightarrow \bigcup_{p\in U}^{}(\{p\}\times T_{p}M) \\
	       & (x, u)\longmapsto (\varphi (x), \varphi '(x)\cdot u)
\end{align*}
É importante notar que TM não é um produto, ou seja, ele não é meramente dado por \(M\times V\) com V um espaço vetorial.

Para que possamos definir o que são as derivadas no contexto puramente de variedades, já temos as ferramentas necessárias! Sejam M e N variedades de \(\mathbb{R}^{k}\) e \(\mathbb{R}^{\ell }\) respectivamente, \(f:M\rightarrow N\) suave e x um ponto de M; já vimos que existem abertos \(W\subseteq \mathbb{R}^{k}\) e \(F:W\rightarrow \mathbb{R}^{\ell }\) suave, com x em W e
\[
	F|_{W\cap M}=f|_{W\cap M}.
\]
\begin{def*}
	Definimos a \textbf{derivada de f em x} como
	\[
		f'(x) = f'(x)|_{T_{x}M}.\; \square
	\]
\end{def*}
Definindo desta forma, \(f'(x)\cdot T_{x}M\) é um subconjunto de \(T_{f(x)}M\), além de fazer \(f'(x)\) bem-definida.
Assim como o usual, a derivada é uma transformação linear e vale a regra da cadeia: se \(f:M\rightarrow N\) e \(g:N\rightarrow P\) são suaves, então
\[
	(g\circ f)'(x)=g'(f(x))\cdot f'(x).
\]
No caso dos difeomorfismos, as dimensões são preservadas! Se f é um difeomorfismo entre M e N, vale \(\mathrm{dim}_{}(M) = \mathrm{dim}_{}(N)\).
\end{document}
